%document class
\documentclass[10pt,oneside]{book}
%%%% Page Info + Commands %%%%%{

%packages
\usepackage{pgfplots}
\pgfplotsset{compat=1.18}
\usepackage{amsmath}
\usepackage{amsfonts}
\usepackage{amsthm,letltxmacro}
\usepackage{amssymb}
\usepackage{enumerate}
\usepackage{enumitem}
\usepackage[dvipsnames]{xcolor}
\usepackage{changepage}

%This will shrink the page
\newcommand{\smallpage}{  \setlength \oddsidemargin{.25in}
            \setlength \textwidth{5.5in}
            \setlength \topmargin{0in}
            \setlength \textheight{9in}}


% Commands for sets of numbers
\newcommand{\Z}{\mathbb{Z}}\newcommand{\R}{\mathbb{R}}\newcommand{\N}{\mathbb{N}}

% gordie spacing and boxing commands
\newcommand{\gspc}{\vspace{0.13cm}} % 0.5 space
\newcommand{\gline}{\gspc\\} % 1.5 space
\newcommand{\gbline}{\gspc\gspc\\} % double space
\newcommand{\gbox}[1]{\fbox{\parbox{\linewidth}{#1}}} % multiline box command
\newcommand{\gmod}{\,\text{mod}\,}


\newcommand{\gimage}[3]{
\begin{figure}[h]   
    \centering          
    \includegraphics[width=#2\textwidth]{#1}  
    \caption{#3}
    \label{fig:#1}
\end{figure}\\
}

\newcommand{\centerit}[1]{{\begin{center} #1 \end{center}}}
\newcommand{\ul}[1]{\underline{#1}}

%%%%%%%% command for graphics %%%%%%%%%%%%%
\usepackage{fancyhdr}
%}

%%%% Page Setup %%%%%%%
\smallpage\sloppy
\pagestyle{fancy}
\lhead{\large{\textbf{Problem Set 8}}}
%\rfoot{\thepage/\pageref{LastPage} }
\setlength{\headheight}{10pt}


% Proof writing
\LetLtxMacro\oldproof\proof
\let\endoldproof\endproof
\renewenvironment{proof}[1][\proofname]
  {\vspace{0.2cm}\begin{adjustwidth}{2em}{2em}
   \oldproof[#1]}
  {\endoldproof
\end{adjustwidth}}

\begin{document}

\newcommand{\cg}[1]{\color{OliveGreen}#1\color{white}}
\newcommand{\cb}[1]{\color{BlueViolet}#1\color{white}}

\subsection*{Section 4.4}
\begin{enumerate}
	\item Solve the following linear congruences:
	\begin{enumerate}
		\item $25x\equiv 15\,\gmod 29$\gline
		First, we find the gcd between 25 and 29:
		\begin{align*}
			29 &= 1(25) + 4\\
			25 &= 6(4) + 1\\
			4 &= 4(1) + 0
		\end{align*}
		Thus our gcd is 1. We know that this problem has a solution because $1\mid 15$. \gline
		Working backwards from the Euclidean algorithm gives:
		\begin{align*}
			1&= 25 - 6(4)\\
			1&= 25 - 6(29 - 25)\\
			1&= 7(25) - 6(29)
		\end{align*}
		This fortunately gives us an $x_0$ of $7\cdot 15 \equiv 18 \gmod 29$. Plugging into our equation gives:
		\begin{align*}
			x &= x_0 +\frac{29}{\gcd(25,29)}t\\
			x &= 18 + 29t\\
			&\equiv 18 \gmod 29
		\end{align*}
		Becaue $\gcd(25, 29) = 1$, we know there is one congruence, and that is 18.
		\item [(d)] $36x\equiv 8 \gmod 102$\gline
		First, we find the gcd between 36 and 102:
		\begin{align*}
			102 &= 2(36) + 30\\
			36 &= 1(30) + 6 \\
			30 &= 5(6) + 0
		\end{align*}
		Thus our gcd is 6. Because $6\nmid 8$, there are no solutions. 
		\item [(e)] $34x \equiv 60 \gmod 98$
		First, we find our gcd:
		\begin{align*}
			98 &= 2(34) + 30 \\
			34 &= 1(30) + 4\\
			30 &= 7(4) + 2\\
			4 &= 2(2) + 1
		\end{align*}
		Thus our gcd is 2. Because $2\mid 60$, we continue with the reverse euclidean:
		\begin{align*}
			2&= 30 - 7(4)\\
			2&= 30 - 7(34 - 30)\\
			2&= 8(30) - 7(34)\\
			2&= 8(98 - 2(34)) - 7(34)\\
			2&= 8(98) - 23(34)
		\end{align*}
		Thus, we know that our $x_0$ is $-23\cdot 30\equiv -4\gmod 98$.
		\pagebreak
		Now, we know that our gcd is 2 so we will have 2 solutions such that:
		\begin{align*}
			x &= x_0 + \frac{n}{\gcd(a, n)}t \\
			&= -4 + \frac{98}{2}t\\
			&= -4 + 49t
		\end{align*}
		Hence, our two solutions are 45 and 94 for $t = 0, 1$. 
	\end{enumerate}
	\item [19.] Obtain the eight incongruent solutions of the linear congruence $3x + 4y \equiv 5\gmod 8$. \gline
	To solve this problem, we will use the anology with Theorem 4.7 provided on page 81 that for $\gcd(a,n)=1$, we can write $ax\equiv c - by(\gmod n)$. In this problem, because $3$ is relatively prime to 5, we can do the same.\gline
	We represent this problem in the form:
	\begin{align*}
		3x \equiv 5 - 4y\gmod 8
	\end{align*}
	Now, we simply solve for an incongruence for each version of $y$ under mod 8:
	\begin{center}
		\begin{tabular}{|c|c|c|c|}\hline
			$y$ & Eq. & Work& $x$ \\\hline
			0 & $3x\equiv 5\gmod 8$ & $3(-1)\equiv 5\gmod 8$ & 7\\\hline
			1 & $3x\equiv 1 \gmod 8$ & $3(3)\equiv 1 \gmod 8$ & 3\\\hline
			2 & $3x\equiv 5 \gmod 8$ & $3(-1)\equiv 3 \gmod 8$ & 7\\\hline
			3 & $3x\equiv 1\gmod 8$ & $3(3)\equiv 5 \gmod 8$ & 3\\\hline
			4 & $3x \equiv 5\gmod 8$ & $3(-1) \equiv 7\gmod 8$ & 7\\\hline
			5 & $3x \equiv 1 \gmod 8$ & $3(3)\equiv 4\gmod 8$ & 3\\\hline
			6 & $3x \equiv 5 \gmod 8$ & $3(-1)\equiv 5\gmod 8$ & 7\\\hline
			7 & $3x \equiv 1\gmod 8$ & $3(3)\equiv 1 \gmod 8$ & 3\\\hline
		\end{tabular}
	\end{center}
	Thus, we end with the incongruent solutions: 
	\[(7,0), (7,2), (7,4), (7,6), (3,1), (3,3), (3,5), (3,7)\]
\end{enumerate}


\end{document}
